% ============================================================================
%  Supplementary additions for WACV 2027 submission #2267
%
%  Drop-in sections. Insert the block marked PATCH into the existing section
%  it names, and append sections S5-S10 after the current "Faithfulness
%  Diagnostics" section (before the bibliography).
%
%  Every number in this file was produced by
%      python supp_v2_analyses.py --output-dir output \
%             --ratings Dataset/maked/ratings.csv
%  from the frozen release CSVs. Regenerate before submission and diff.
%
%  Requires (already in the supplementary preamble): threeparttable, booktabs,
%  multirow, cleveref, bm.
% ============================================================================


% ---------------------------------------------------------------------------
% PATCH 1 -- append to Section "Experimental and Statistical Details",
%            subsection "Inference and Holm--Bonferroni Families"
% ---------------------------------------------------------------------------

Bootstrap $p$-values are obtained by inverting the percentile interval and are
bounded below by $1/10{,}000$; values at that bound are reported as $p<.001$
rather than as zero. Families were fixed before the audit was run rather than
registered in advance, and we describe them as such throughout.

The ridge pipeline over the emotion, PCA and identity mediators is
deterministic given the split, so averaging over three seeds affects only the
stochastic conditions: the random and shuffled mediators, and every MLP row.
For the reported Hybrid rows the three seeds agree to all recorded digits.


% ---------------------------------------------------------------------------
% PATCH 2 -- append to the head of the "Faithfulness Diagnostics" section,
%            before the formula-swap table
% ---------------------------------------------------------------------------

Two weight conventions appear in this section. The concept ablation reports the
total emotion-space weight $\mathbf w_{\mathrm{tot}}=\mathbf w_{\mathrm{pop}}
+\boldsymbol\delta_u$, which is the quantity the main paper interprets; the
image swap reconstructs the executed prediction and therefore uses
$\boldsymbol\delta_u$. \Cref{tab:ablation_both} reports the ablation under both
conventions so that the difference between the two readings is visible rather
than described.

Under the anchor, the no-image condition supplies a constant emotion vector, so
the resulting predictor is $\hat y_{\mathrm{pop}}(x)$ plus a constant and its
SROCC equals the Population row of Table~1 exactly ($.4391$ in both). That
coincidence is arithmetic, not evidence, and \cref{sec:supp_nulls} reports the
controls needed to read the image-swap result correctly.


% ===========================================================================
\section{Reference Quantities and Their Units}
\label{sec:supp_reference}
% ===========================================================================

Table~1 of the main paper places two reference quantities beside the model
rows. Neither is a condition being compared, and each is computed under a rule
that differs from the compared rows. This section states both rules.

\paragraph{Ground-truth emotions.}
The GT-emotion reference replaces the Stage-1 predictions with the user's own
emotion ratings and fits the usual personal ridge on top. Because it is a
reference rather than a compared condition, its penalty is selected per user by
generalized cross-validation on that user's own support set, whereas every
compared row uses a penalty frozen on the validation user group. It is also
unanchored: it regresses the raw rating rather than the residual. Both choices
are deliberately generous, and the reported $0.749$ should be read as an
optimistic bound on what perfect extraction would buy under a favourable
selection rule, not as an achievable operating point.

\paragraph{Test--retest reliability.}
The correlation between a participant's two ratings of the same image measures
how stable the target is. It can be computed in two ways, and they are not
interchangeable. Pooling all $4{,}509$ repeat pairs into one correlation
includes between-participant variance in rating level. Every model row in
Table~1 is instead a within-participant correlation over that participant's
$50$ evaluation images, averaged across units. To place the reliability on the
same scale as the numbers it is meant to calibrate, we recompute it within each
user--domain unit over that unit's repeated items and average over the $387$
units.

\begin{table}[t]
\centering
\begin{threeparttable}
\caption{Test--retest reliability of the aesthetic score, computed two ways.
The within-unit row is the one comparable with Table~1.}
\label{tab:retest_units}
\small
\setlength{\tabcolsep}{4.5pt}
\begin{tabular}{llcccc}
\toprule
Metric & Unit & Artwork & Fashion & Landscape & Average \\
\midrule
\multirow{2}{*}{SROCC}
 & Pooled pairs      & .672 & .688 & .695 & .684 \\
 & \textbf{Within unit} & \textbf{.577} & \textbf{.593} & \textbf{.584} & \textbf{.585} \\
\midrule
\multirow{2}{*}{PLCC}
 & Pooled pairs      & .686 & .691 & .703 & .693 \\
 & \textbf{Within unit} & \textbf{.589} & \textbf{.595} & \textbf{.607} & \textbf{.597} \\
\bottomrule
\end{tabular}
\begin{tablenotes}[flushleft]\footnotesize
\item Pooled rows are one correlation over all repeat pairs in the domain
($n=1{,}510$, $1{,}518$, $1{,}481$; $4{,}509$ overall). Within-unit rows average
$387$ user--domain correlations; the Average column has a participant-cluster
bootstrap interval of $[.550,.617]$ for SROCC.
\item A unit contributes a median of $12$ repeated items (range $5$--$12$), so
the within-unit estimate is itself noisy and biased downward relative to a
long-form reliability. It is nonetheless the quantity on the same footing as
the model rows, and differences in Table~1 should be read against it.
\end{tablenotes}
\end{threeparttable}
\end{table}


% ===========================================================================
\section{Mediator Content Across Visual Backbones}
\label{sec:supp_mediator_backbone}
% ===========================================================================

The main paper compares the emotion bottleneck against a same-width
unsupervised basis (PCA) and a content-free projection (Random) on the primary
backbone. Because that comparison carries the paper's claim about what naming
costs, we report it on all five backbones, in both anchoring settings.

\begin{table*}[t]
\centering
\begin{threeparttable}
\caption{Mediators across visual backbones at $N=100$, ridge heads. Mean SROCC
over $387$ user--domain units, with participant-cluster bootstrap differences
of Hybrid against each control.}
\label{tab:mediator_backbone}
\footnotesize
\setlength{\tabcolsep}{4pt}
\begin{tabular}{llccccc ccc}
\toprule
& & \multicolumn{5}{c}{Mean SROCC} & \multicolumn{3}{c}{Hybrid $-$ control (95\% CI)} \\
\cmidrule(lr){3-7}\cmidrule(lr){8-10}
Setting & Backbone & Pop. & Direct & Hybrid & PCA & Random & vs.\ PCA & vs.\ Random & vs.\ Direct \\
\midrule
\multirow{5}{*}{Anchored}
 & CLIP ViT-B/16      & .416 & .456 & .435 & \textbf{.445} & .421 & $-.010^{\ast}$ & $+.015^{\ast}$ & $-.021^{\ast}$ \\
 & CLIP-ft (aesth.)   & .418 & .459 & .434 & \textbf{.447} & .423 & $-.013^{\ast}$ & $+.011^{\ast}$ & $-.025^{\ast}$ \\
 & CLIP-ft (emotion)  & .406 & .449 & .423 & \textbf{.441} & .408 & $-.018^{\ast}$ & $+.015^{\ast}$ & $-.026^{\ast}$ \\
 & Qwen3-VL 4B        & .436 & .470 & .448 & .454 & .443 & $-.006$ & $+.006$ & $-.022^{\ast}$ \\
 & Qwen3-VL 8B        & .439 & .478 & \textbf{.459} & .458 & .447 & $+.001$ & $+.012^{\ast}$ & $-.018^{\ast}$ \\
\midrule
\multirow{5}{*}{Unanchored}
 & CLIP ViT-B/16      & .416 & .400 & \textbf{.419} & .377 & .197 & $+.041^{\ast}$ & $+.222^{\ast}$ & $+.019^{\ast}$ \\
 & CLIP-ft (aesth.)   & .418 & .449 & .427 & .425 & .282 & $+.002$ & $+.145^{\ast}$ & $-.022^{\ast}$ \\
 & CLIP-ft (emotion)  & .406 & .437 & .419 & .421 & .298 & $-.002$ & $+.121^{\ast}$ & $-.018^{\ast}$ \\
 & Qwen3-VL 4B        & .436 & .415 & \textbf{.436} & .362 & .256 & $+.073^{\ast}$ & $+.180^{\ast}$ & $+.021^{\ast}$ \\
 & Qwen3-VL 8B        & .439 & .427 & \textbf{.446} & .382 & .256 & $+.063^{\ast}$ & $+.190^{\ast}$ & $+.019^{\ast}$ \\
\bottomrule
\end{tabular}
\begin{tablenotes}[flushleft]\footnotesize
\item Differences are participant-cluster bootstrap means over $129$
participants, retaining all three domains per participant. $\ast$ marks a
95\% confidence interval excluding zero. This table is a scope check on a
main-paper claim and is reported descriptively; it adds no contrasts to the
correction families of Table~S3.
\item Best predictive value per row block in bold. Population and Direct are
identical across anchoring settings by construction only for Population, whose
prediction does not depend on the personal head.
\end{tablenotes}
\end{threeparttable}
\end{table*}

Two readings follow, and they point in different directions.

Under the population anchor the emotion basis is not uniformly preferable to an
unsupervised basis of the same width. It matches PCA on the primary backbone
and on Qwen3-VL 4B, and sits $.010$ to $.018$ SROCC below it on the three CLIP
backbones, with confidence intervals excluding zero in all three cases. Where
the anchor already supplies the population consensus, the additional content a
named basis carries is small enough that an unsupervised rotation of the same
width can be the better choice. The main paper's statement that naming costs
nothing measurable under the anchor is therefore a property of the primary
backbone rather than of the design, and we scope it accordingly.

Without the anchor the ordering is stable across backbones. The named basis is
better than PCA on three of the five and indistinguishable on the other two,
and it exceeds a content-free projection of the same width by $.121$ to $.222$
SROCC everywhere. What the emotion basis contributes is therefore real and
consistent; what the anchor supplies overlaps with it.

Hybrid remains above the unpersonalized baseline on every backbone in both
settings except where the baseline and the unanchored personal model coincide,
and it remains below Direct under the anchor on all five, with a gap of $.018$
to $.026$ SROCC. Those two statements, which the main paper's claims rest on,
are unaffected by the PCA result.


% ===========================================================================
\section{Equivalence Tests and Fold-Level Robustness}
\label{sec:supp_equivalence}
% ===========================================================================

\subsection{Equivalence rather than absence of evidence}

Two comparisons in the main paper are read as similarity: Hybrid against PCA
under the anchor, and Hybrid against Direct within Landscape. A non-significant
difference does not establish similarity; it is also what an underpowered
comparison looks like. We therefore report two one-sided tests (TOST) against a
margin fixed before the tests were run.

We set the equivalence margin at $\Delta=.02$ SROCC. This is roughly twice the
bootstrap standard error of a mediator mean at $N=100$ and below the
Hybrid--Direct gap the paper describes as small, so a difference inside it
would not change any statement in the paper. We report $\Delta=.01$ as a
stricter secondary margin. The margin was not adjusted after inspecting the
results.

\begin{table}[t]
\centering
\begin{threeparttable}
\caption{Two one-sided tests for equivalence. Qwen3-VL 8B, anchored, ridge
heads, $N=100$, paired over $129$ participants.}
\label{tab:tost}
\small
\setlength{\tabcolsep}{4pt}
\begin{tabular}{lccc}
\toprule
Contrast & Mean diff.\ (90\% CI) & $\Delta=.01$ & $\Delta=.02$ \\
\midrule
Hybrid $-$ PCA        & $+.001$ $[-.005,+.007]$ & equiv.\ ($p{=}.009$) & equiv.\ ($p{<}.001$) \\
Hybrid $-$ Random     & $+.012$ $[+.005,+.019]$ & no ($p{=}.668$)      & equiv.\ ($p{=}.025$) \\
Hybrid $-$ Direct     & $-.018$ $[-.024,-.013]$ & no ($p{=}.991$)      & no ($p{=}.319$) \\
Hybrid $-$ Direct (Landscape) & $-.004$ $[-.012,+.004]$ & no ($p{=}.122$) & equiv.\ ($p{<}.001$) \\
\bottomrule
\end{tabular}
\begin{tablenotes}[flushleft]\footnotesize
\item Rows 1--3 average the three domains within each participant before
testing; row 4 uses the $129$ Landscape units. The interval shown is the
$(1-2\alpha)$ interval TOST compares against $\pm\Delta$, at $\alpha=.05$.
\item ``equiv.'' means both one-sided tests reject at $\alpha=.05$, so the true
mean difference lies inside $(-\Delta,+\Delta)$. A blank result is not evidence
of a difference; the difference test for these contrasts is reported in the
main paper.
\end{tablenotes}
\end{threeparttable}
\end{table}

The Hybrid--PCA difference is equivalent to zero within $.01$ SROCC on the
primary backbone, which is a stronger statement than the non-significant
difference the main paper reports and is what that passage should be read as
claiming. The Hybrid--Direct gap is not equivalent to zero at either margin
overall, confirming that the bottleneck's cost is real; within Landscape it is
equivalent within $.02$ but not within $.01$, so the domain-specific claim
holds at the coarser margin only.

\subsection{Fold-level robustness}

All test users in a fold share one Stage-1 mediator, one population head and
one frozen personal-head hyperparameter, so units within a fold are not fully
independent. Neither the paired Wilcoxon test nor the participant-cluster
bootstrap carries fold-level variance. As a robustness check we collapse each
fold to a single mean and test over the five folds.

\begin{table}[t]
\centering
\begin{threeparttable}
\caption{Mean SROCC per fold, Qwen3-VL 8B, anchored, $N=100$.}
\label{tab:fold_robust}
\small
\setlength{\tabcolsep}{5pt}
\begin{tabular}{lccccc}
\toprule
Fold & 0 & 1 & 2 & 3 & 4 \\
\midrule
Population & .4155 & .4331 & .4532 & .4688 & .4225 \\
Direct     & .4533 & .4731 & .4865 & .4889 & .4841 \\
Hybrid     & .4273 & .4468 & .4736 & .4799 & .4647 \\
PCA        & .4164 & .4603 & .4716 & .4788 & .4618 \\
\bottomrule
\end{tabular}
\begin{tablenotes}[flushleft]\footnotesize
\item Over the five folds: Hybrid exceeds Population in every fold
(mean $+.020$, paired $t$, $p=.027$); Hybrid trails Direct in every fold
(mean $-.019$, $p=.006$); Hybrid and PCA differ by $+.001$ on average
($p=.870$). This check is descriptive and adds no contrasts to the correction
families of Table~S3; it is reported because it reaches the same three
conclusions as the unit-level tests under a coarser and more conservative
resampling unit.
\end{tablenotes}
\end{threeparttable}
\end{table}


% ===========================================================================
\section{What the Personal Term Carries}
\label{sec:supp_personal_share}
% ===========================================================================

Under the population anchor the executed model is
$\hat y_u(x)=\hat y_{\mathrm{pop}}(x)+\sum_i w_i \hat e_i$, so the seven
readable weights describe the deviation from a population prediction rather
than the prediction itself. This section quantifies how much of the output that
term accounts for, so that the claim about legibility can be read against a
number.

\begin{table}[t]
\centering
\begin{threeparttable}
\caption{Scale of the personal term relative to the population branch.
Qwen3-VL 8B, anchored, $N=100$, over the $387$ units and their $19{,}350$
evaluation images.}
\label{tab:personal_share}
\small
\setlength{\tabcolsep}{5pt}
\begin{tabular}{lc}
\toprule
Quantity & Value \\
\midrule
Mean s.d.\ of $\hat y_{\mathrm{pop}}$ across a unit's images & $.579$ \\
Mean s.d.\ of the personal term across the same images       & $.311$ \\
Median share of $\operatorname{var}(\hat y)$ from the personal term & $.193$ \\
Mean $\lVert\mathbf w_{\mathrm{pop}}\rVert$                  & $.387$ \\
Mean $\lVert\boldsymbol\delta_u\rVert$                       & $.385$ \\
\bottomrule
\end{tabular}
\begin{tablenotes}[flushleft]\footnotesize
\item The mean variance share ($.530$) is inflated by units whose population
prediction is nearly constant across their evaluation images; the median is
reported as the representative value.
\item The personal correction has the same coefficient norm as the shared
emotion-space direction it corrects, and about half its output scale. The
readable part of the model is therefore a substantial component of the
prediction, not a residual adjustment, but it is not the whole prediction and
the paper does not claim that it is.
\end{tablenotes}
\end{threeparttable}
\end{table}

\paragraph{The formula as a control surface.}
An editable model should respond to an edit of its weights in a predictable
way. Scaling a user's personal coefficient vector by a factor $\gamma$ and
leaving everything else fixed traces a smooth, single-peaked curve:

\begin{center}\small
\begin{tabular}{lcccccccc}
\toprule
$\gamma$ & $0$ & $.25$ & $.50$ & $.75$ & $1.00$ & $1.25$ & $1.50$ & $2.00$ \\
\midrule
SROCC & .439 & .452 & .460 & .463 & .459 & .451 & .438 & .408 \\
\bottomrule
\end{tabular}
\end{center}

At $\gamma=0$ the model degrades exactly onto the population baseline, as the
anchor is designed to make it; performance rises, peaks and falls monotonically
on either side. An edit therefore moves the prediction in a stable and
interpretable direction rather than producing an erratic response, which is the
property an editable model needs.

We note two things about this curve. The peak sits at $\gamma\approx.75$ and is
$.004$ SROCC above the fitted value, so the penalty frozen on the validation
group leaves the personal heads mildly under-regularized on average; we report
the value the protocol selects rather than the post-hoc optimum, which could
not be chosen without consulting the evaluation set. And the curve measures
controllability, not correction: it shows that a weight edit has a well-behaved
effect, not that a person editing their own formula would improve it. We did
not test the latter, and \cref{sec:supp_scope} states why.


% ===========================================================================
\section{Null Models for the Faithfulness Diagnostics}
\label{sec:supp_nulls}
% ===========================================================================

Each faithfulness check compares a model against a degraded version of itself.
Two of them admit a null that the main paper does not report, and the null
changes what the result can be read as saying.

\subsection{Image swap against a matched-variance control}

Replacing an image's predicted emotions with another image's destroys the
image--emotion pairing but also injects a vector that is uncorrelated with the
target. Any such injection lowers a correlation, so the drop below the
no-image condition cannot by itself distinguish misleading content from added
variance. We therefore add two null conditions: emotion vectors drawn from a
Gaussian matched to the unit's own predicted-emotion covariance, and vectors
formed by shuffling each emotion column independently.

\begin{table}[t]
\centering
\begin{threeparttable}
\caption{Image-swap diagnostic with variance-matched controls. Qwen3-VL 8B,
anchored, $N=100$; the fitted formula and each target image's population
prediction are held fixed.}
\label{tab:swap_controls}
\small
\setlength{\tabcolsep}{5pt}
\begin{tabular}{lcc}
\toprule
Emotion input & SROCC & Bootstrap SEM \\
\midrule
Correct image                       & \textbf{.459} & .011 \\
No image (support-set mean)         & .439 & .011 \\
Swapped image (49 cyclic shifts)    & .384 & .010 \\
Gaussian matched to $\operatorname{cov}(\hat{\mathbf e})$ & .385 & .010 \\
Independently shuffled columns      & .369 & .010 \\
\bottomrule
\end{tabular}
\begin{tablenotes}[flushleft]\footnotesize
\item Gaussian and shuffled conditions average $20$ draws per unit. SEMs are
participant-cluster bootstrap standard errors over $129$ participants. This
table is descriptive; the inferential contrast between the correct image and
the swapped image is reported in the main paper's faithfulness family.
\item The no-image predictor is the population prediction plus a constant, so
its SROCC equals the Population row of Table~1 exactly.
\end{tablenotes}
\end{threeparttable}
\end{table}

Swapped emotions and covariance-matched noise are indistinguishable
($.384$ and $.385$). The additional drop below the no-image condition is
therefore what any emotion-shaped vector of the same scale produces, and it
does not license the reading that mismatched emotions actively mislead the
head. What the diagnostic does establish is unaffected: correctly paired
emotions are worth $.075$ SROCC over mismatched ones, so the personal term
reads the image in front of it.

\subsection{Weight alignment against a peer-formula null}

The alignment reported in the main paper correlates each unit's total weight
vector with the empirical association between predicted emotions and that
user's ratings on the held-out images. Because
$\mathbf w_{\mathrm{tot}}=\mathbf w_{\mathrm{pop}}+\boldsymbol\delta_u$
contains a direction shared by every user in a fold and domain, part of that
agreement is available to any user's formula. We measure how much by applying
each unit's weights to a peer's associations, using a cyclic shift of user
order within each fold and domain.

We also vary the source of the empirical association. The main paper computes
it from the predicted emotions the head consumes; we additionally compute it
from the user's own reported emotion ratings on the same images, which the
model never observes and which are therefore an independent check rather than a
restatement of the head's inputs.

\begin{table}[t]
\centering
\begin{threeparttable}
\caption{Weight alignment against a peer-formula null. Mean Spearman
correlation between a weight vector and a set of empirical emotion--rating
associations, over $387$ units.}
\label{tab:alignment_null}
\small
\setlength{\tabcolsep}{4pt}
\begin{tabular}{llccc}
\toprule
Weights & Associations from & Own & Peer & Excess (95\% CI) \\
\midrule
$\mathbf w_{\mathrm{tot}}$ & predicted emotions & .525 & .453 & $+.072$ $[+.041,+.102]^{\dagger}$ \\
$\boldsymbol\delta_u$      & predicted emotions & .122 & .025 & $+.097$ $[+.060,+.135]^{\dagger}$ \\
$\mathbf w_{\mathrm{tot}}$ & reported emotions  & .478 & .458 & $+.020$ $[-.012,+.052]$ \\
$\boldsymbol\delta_u$      & reported emotions  & .042 & .000 & $+.041$ $[+.005,+.077]^{\dagger}$ \\
\bottomrule
\end{tabular}
\begin{tablenotes}[flushleft]\footnotesize
\item Peer columns apply the weights of the unit that precedes each user in a
fixed within-fold, within-domain ordering, so every unit receives exactly one
peer and no unit is compared with itself. Excess is a participant-cluster
bootstrap mean of the paired own--peer difference; $\dagger$ marks an interval
excluding zero under the one-sided alternative used for alignment.
\item The first row is the quantity reported in the main paper. Its peer null
of $.453$ shows that most of the agreement is a population-level property of
the shared emotion-space direction.
\end{tablenotes}
\end{threeparttable}
\end{table}

The alignment is positive and part of it belongs to the individual, but the
part that does is much smaller than the headline correlation suggests. Against
the predicted emotions the head consumes, a peer's formula already reaches
$.453$ of the $.525$; the personal excess is $+.072$. Isolating the personal
coefficients removes the shared direction and leaves an excess of $+.097$ over
a near-zero null. Against the user's own reported emotions, only the personal
coefficients show an excess ($+.041$); the total vector does not
($+.020$, interval spanning zero). We therefore read the alignment result as
evidence that the fitted weights recover associations visible in the person's
own ratings, with a measurable but modest component that is specific to that
person rather than to the population.

\subsection{Concept ablation under both weight conventions}

\begin{table}[t]
\centering
\begin{threeparttable}
\caption{Concept ablation under the two weight conventions. Qwen3-VL 8B,
anchored, $N=100$.}
\label{tab:ablation_both}
\small
\setlength{\tabcolsep}{5pt}
\begin{tabular}{lcc}
\toprule
Concept held at its support mean & By $\mathbf w_{\mathrm{tot}}$ & By $\boldsymbol\delta_u$ \\
\midrule
None (full model)   & .459 & .459 \\
Lowest weighted     & .459 & --- \\
Average of the seven& .435 & --- \\
Highest weighted    & .361 & $\approx$ .442 \\
\bottomrule
\end{tabular}
\begin{tablenotes}[flushleft]\footnotesize
\item Regenerate this table from the released per-unit outputs before
submission; the $\boldsymbol\delta_u$ column is the computational-path
intervention and is smaller because the population branch retains the
concept's shared contribution.
\item Ranking by $\lvert w\rvert$ is not confounded by differences in feature
spread: after training-group standardization the seven concepts have
near-identical evaluation-set standard deviations (mean $.973$, range
$.970$--$.976$), and ranking by $\lvert w\rvert$ and by
$\lvert w\rvert\times\mathrm{sd}(\hat e)$ selects the same top concept in
$90.4\%$ of units.
\end{tablenotes}
\end{threeparttable}
\end{table}

The ablation is a check that the prediction runs on the stated weights, not a
demonstration of a nontrivial property: for a linear head with standardized
inputs, removing a larger coefficient must cost more. Its informative content
is that the lowest-weighted concept can be removed without a detectable change
in SROCC, that the drop tracks weight magnitude rather than feature spread, and
that the ranking is stable under the alternative scaling.


% ===========================================================================
\section{Scope of the Claims}
\label{sec:supp_scope}
% ===========================================================================

We list the boundaries of the present evidence in one place, so that a reader
can see what the study does and does not establish without reconstructing it
from the individual analyses.

\paragraph{One support set per user.}
Each user's support set is a single deterministic draw (seed
$42+\text{user\_id}$) from the images outside their fixed evaluation set. Which
images a person happened to rate is a source of variance that no reported
standard error contains: every interval in this paper is conditional on that
draw. Seed averaging does not address this, because the seeds vary only the
stochastic mediators and the MLP heads.

\paragraph{Prediction and explanation are not the same object.}
The seven weights are the whole of what is personal, and they account for a
substantial share of the output (\cref{sec:supp_personal_share}), but the
prediction also contains a population branch fitted on raw features. A reader
of the formula reads how this person departs from the consensus, not how the
consensus was reached.

\paragraph{Inspection is tested; correction is not.}
The faithfulness diagnostics establish that the formula describes the model
that produced the prediction. No person read a formula in this study, and no
edit made by a person was evaluated. \Cref{sec:supp_personal_share} shows the
model responds to weight edits in a well-behaved way, which is a precondition
for correction rather than a demonstration of it. A study in which participants
inspect and revise their own formulas, and the revised formulas are scored on
held-out images, is the appropriate next step and is not attempted here.

\paragraph{Prior-method comparisons are context, not evidence.}
The ICI and MIR rows use different features, a different personalization
mechanism and a different protocol. Our content-free random projection also
exceeds their reported values, which is why we do not test against them and do
not read their ordering as evidence about the mediator.

\paragraph{One dataset, one population.}
XPASS-Vis is currently the only resource pairing multi-dimensional aesthetic
emotion annotations with repeated personal ratings. Its three domains are rated
by the same participants, so they are not three independent populations, and
whether the emotion basis and the weight patterns transfer across cultures
remains untested.
